\documentclass[12pt,a4paper,oneside]{article}
\usepackage[brazil]{babel}
\usepackage[utf8]{inputenc}
\usepackage{olymp}

\setcounter{problem}{2}

\begin{document}

\begin{problem}{Início das aulas}{data}{2}

As aulas de INF110 começaram no dia 20 de março. A maioria dos estudantes se matriculou na disciplina antes das aulas começarem, mas alguns se matricularam depois, porque estavam aguardando vaga ou as chamadas do SISU.

Nesta questão você deve ler a data de matrícula de um estudante e informar se sua matrícula foi feita antes, no dia, ou depois do início das aulas.

%\Notes
%
%Observações, se tiver.

\InputFile

A entrada é composta por dois valores inteiros $D$ e $M$, que indicam respectivamente o dia e o mês da matrícula de um estudante. Restrição: as datas serão sempre dias válidos.


%\Notes

\OutputFile

Seu programa deve gerar apenas uma linha na saída, contendo um dos textos a seguir:
\begin{itemize}
    \item \texttt{antes}: se o dia da matrícula foi antes de 20/3
    \item \texttt{no dia}: se o dia da matrícula foi exatamente 20/3
    \item \texttt{depois}: se o dia da matrícula foi depois de 20/3
\end{itemize}

\Example

\begin{example}
\exmp{
19 3
}{
antes
}%
\end{example}

\begin{example}
\exmp{
19 4
}{
depois
}%
\end{example}

\begin{example}
\exmp{
20 3
}{
no dia
}%
\end{example}

\begin{example}
\exmp{
20 2
}{
antes
}%
\end{example}

%comment

\end{problem}

\end{document}
